\documentclass[12pt,a4paper]{article}

\usepackage[utf8]{inputenc}
\usepackage[T1]{fontenc}
\usepackage[
  a4paper,
  top=2.5cm, bottom=3cm,
  left=2.5cm, right=3.8cm,
  marginparwidth=2.8cm,
  marginparsep=0.5cm
]{geometry}
\usepackage{xcolor}
\usepackage{marginnote}
\usepackage{microtype}
\usepackage{parskip}
\usepackage{titlesec}
\usepackage{float}
\usepackage{amsmath}
\usepackage{booktabs}
\usepackage{multirow}
\usepackage{caption}

% ── Colours ───────────────────────────────────────────────────────────────────
\definecolor{respblue}{RGB}{0,48,128}
\definecolor{cidred}{RGB}{204,0,0}

% ── Custom commands ───────────────────────────────────────────────────────────
% Comment identifier  (red, underlined)
\newcommand{\cid}[1]{{\color{cidred}\underline{#1}}}

% Margin section reference
\newcommand{\mref}[1]{%
  \marginnote{%
    \color{respblue}\scriptsize
    \rule[-0.3ex]{0.6pt}{1.4ex}\hspace{4pt}%
    \parbox[t]{2.4cm}{\raggedright #1}%
  }%
}

% Response environment (blue text)
\newenvironment{response}{%
  \par\color{respblue}%
}{\par}

% Block-quote of text added to the proposal
\newenvironment{quotedaddition}{%
  \begin{quote}\color{respblue}%
}{\end{quote}}

% ── Section heading style ─────────────────────────────────────────────────────
\titleformat{\section}[block]
  {\normalfont\large\bfseries}{}{0pt}{}
  [\vspace{-6pt}\rule{\linewidth}{0.8pt}]
\titlespacing{\section}{0pt}{18pt}{8pt}

\setlength{\parskip}{0.6em}
\setlength{\parindent}{0pt}

% ══════════════════════════════════════════════════════════════════════════════
\begin{document}

\begin{center}
  {\large\bfseries Responses to the Examination Committee}\\[5pt]
  {\normalsize Development of a Predictive Model for Diabetic Foot Risk Using\\
  Podoscope Data with XAI-Based Explanations}\\[6pt]
  {\normalsize July 6, 2026}
\end{center}

\medskip\hrule\medskip

\setlength{\parindent}{1.5em}

We have carefully read the comments of the Examination Committee regarding our thesis proposal and have extensively revised the proposal to address the points raised during the examination on July~1, 2026.

We would like to thank Asst.\ Prof.\ Saneh Khuenkaew and Assoc.\ Prof.\ Dr.\ Worapan Kusakunniran for their thorough and constructive feedback. We believe the manuscript has improved greatly as a result of these comments. We carefully addressed all eleven points to improve the proposal and believe the modifications have addressed the concerns raised.

With this letter, we provide each comment from the examination, followed by our response explaining the reasoning behind the concern and describing how the proposal has been revised. The annotations in the margin indicate the section of the proposal that was modified or extended.

\setlength{\parindent}{0pt}

% ══════════════════════════════════════════════════════════════════════════════
\section*{Response to Asst.\ Prof.\ Saneh Khuenkaew}

% ── C1 ────────────────────────────────────────────────────────────────────────
\cid{C1} The inclusion criteria should specify Type~2 Diabetes Mellitus patients only, as Type~1 and Type~2 DM have different pathogeneses. Congenital foot deformities belong in the inclusion criteria rather than the exclusion criteria, since they are indicators of IWGDF risk category~2 or above. The criteria should also state explicitly that participants must be cognitively alert and able to follow verbal commands.

\mref{\S~3.2.5\\[2pt]Inclusion and\\Exclusion\\Criteria}
\begin{response}
\underline{Response:} We thank the Committee for this guidance and agree with all three points. On the type of diabetes, Type~1 and Type~2 DM have fundamentally distinct pathogeneses, with Type~1 arising from autoimmune destruction of insulin-producing pancreatic beta cells and Type~2 from progressive insulin resistance, which may produce different complication patterns relevant to the diabetic foot ulcer risk pathway. Restricting the inclusion criteria to Type~2 DM therefore reduces confounding and improves sample homogeneity.

Regarding foot deformities, we agree with the Committee. Congenital foot deformities produce the same mechanical pressure concentration patterns that elevate ulceration risk, and excluding them on the basis of origin was inconsistent with the study's screening objective. Foot deformity has accordingly been removed from the exclusion criteria, so that patients presenting with a deformity are eligible rather than screened out. Because a deformity indicates IWGDF risk category~2 or above rather than being a prerequisite for enrolment, it is recorded during the foot inspection and incorporated into the risk category assigned to each patient.

Finally, we fully agree that cognitive alertness is a necessary inclusion criterion. The ability to perceive and follow verbal instructions affects the reliability of LOPS assessment, which relies primarily on patient-reported sensation, and is equally necessary for maintaining a correct standing posture during photo-podoscope image acquisition. A patient who cannot follow instructions may compromise both image quality and patient safety during the standing examination. Section~3.2.5 has been updated to reflect all three changes.
\end{response}

% ── C2 ────────────────────────────────────────────────────────────────────────
\cid{C2} The exclusion criterion on lower-limb amputation should be more specific. Only patients with bilateral lower-limb amputation should be excluded, since a patient who still has one weight-bearing limb is able to complete the assessment.

\mref{\S~3.2.5\\[2pt]Exclusion\\Criteria}
\begin{response}
\underline{Response:} We agree that the wording requires refinement. The original wording, ``history of lower-limb amputation'', without further qualification, was broader than intended. It would exclude any patient who has lost a limb, including those who have lost only one leg but retain a functioning weight-bearing foot on the other side. The actual intent of this criterion is simply to exclude patients who cannot stand on the podoscope platform, and a patient with unilateral amputation is still able to do so. Furthermore, patients with a history of amputation fall under IWGDF Category~3, the highest risk stratum, meaning that data from their remaining foot is among the most clinically important the study can collect. Excluding these patients would therefore remove precisely the cases the model most needs to learn from. The exclusion criterion has accordingly been revised to specify ``bilateral lower-limb amputation'', which is the condition that actually prevents a patient from completing the assessment. Section~3.2.5 has been updated accordingly.
\end{response}

% ── C3 ────────────────────────────────────────────────────────────────────────
\cid{C3} The study setting description is not specific enough. Rather than naming the hospital alone, it should identify the diabetic clinic or outpatient diabetic care unit, as this more accurately reflects where the target patient population is seen.

\mref{\S~3.2.3\\[2pt]Data Source}
\begin{response}
\underline{Response:} We agree with this recommendation because we originally mentioned only the hospital name, which did not clearly tell the reader where our study would take place. The target population that we will collect data from is patients with Type~2 Diabetes Mellitus who are followed in the diabetic clinic, where blood glucose monitoring and foot screening are already part of routine care. Therefore, clearly stating ``diabetic clinic'' in Section~3.2.3 better indicates the actual study setting than mentioning only the hospital name.
\end{response}

% ── C4 ────────────────────────────────────────────────────────────────────────
\cid{C4} The 10-gram monofilament test and the Ankle-Brachial Index are the standard clinical tools for diagnosing LOPS and PAD respectively, yet they are not mentioned in the Chapter~2 literature review or in the overall study design. These should be included, as they are the instruments used to assign each patient's IWGDF risk category.

\mref{\S~2.3\\[2pt]Clinical\\Assessment\\Methods\\[4pt]\S~3.3\\[2pt]Data\\Annotation}
\begin{response}
\underline{Response:} We thank the Committee for raising this point. The monofilament test and the ABI are not just background information, because they are the real tools we use to determine each patient's IWGDF risk category. We understand that the reader may think we are only mentioning these tools without actually using them, so we want to make it clear that we have studied them in detail and that we do use them following standard clinical practice. These tools are already described in Section~2.3, but the explanation there is still not detailed enough. We will also make it clear again in Section~3.3 that they are part of the data annotation process.
\end{response}

% ══════════════════════════════════════════════════════════════════════════════
\section*{Response to Assoc.\ Prof.\ Dr.\ Worapan Kusakunniran}

% ── C5 ────────────────────────────────────────────────────────────────────────
\cid{C5} Rather than limiting the model to binary classification, this study should explore whether it can distinguish patients across all four IWGDF risk categories (0, 1, 2, and~3). Demonstrating this staging capability would substantially strengthen the clinical contribution of the work.

\mref{\S~1.3\\Research\\Question (RQ5)\\[4pt]\S~3.7\\Evaluation\\Plan}
\begin{response}
\underline{Response:} Thank you for this interesting suggestion. We agree that it would make the study stronger, but in this work we want to keep the scope limited to binary classification, since our goal is only to show whether a patient should be referred or not. This is not as broad as building a model that can distinguish all four levels, and the result of a multi-class model would be more useful for patient triage and resource allocation than a binary outcome. It may also help show whether plantar pressure images contain enough signal for four-class classification, since conditions such as LOPS and PAD without foot deformity may not leave a clearly visible pressure pattern. For this reason, we regard this as future work rather than part of the current study, because our research focus is on developing a screening model for patients only.
\end{response}

% ── C6 ────────────────────────────────────────────────────────────────────────
\cid{C6} The use of the term ``early detection'' appears to overclaim what this study can actually demonstrate. IWGDF Category~1 is defined by LOPS or PAD, which are neurological and vascular conditions that are not directly observable in plantar pressure images. A patient with Category~1 risk arising from LOPS alone but with no structural foot abnormality may produce a completely normal-looking pressure image.

\mref{\S~3.2.1\\Classification\\Labels}
\begin{response}
\underline{Response:} This concern is well-founded. The IWGDF classification assigns risk on the basis of LOPS, PAD, and structural foot deformity or previous ulcer. A plantar pressure image encodes the mechanical load pattern across the foot surface, which reflects structural changes such as high forefoot loading or deformity-related contact patterns, but it does not capture neurological function or vascular status. A category~1 patient whose risk arises from LOPS or PAD alone, with no structural deformity, may produce a completely normal-looking pressure image, because the relevant clinical signal simply does not appear in this modality. Calling this ``early detection'' implies the model can find risk before it is clinically visible, which is more than a pressure-based image can offer. The more honest framing is ``risk screening'', which positions the tool as a way to flag abnormal pressure patterns for clinical follow-up rather than a replacement for monofilament and ABI testing. The term has been revised throughout the proposal, and Section~3.2.1 has been updated accordingly.
\end{response}

% ── C7 ────────────────────────────────────────────────────────────────────────
\cid{C7} It would be worth investigating whether a single model trained on images from both feet performs adequately, or whether separate models for the left and right foot would be more appropriate.

\mref{\S~1.3\\Research\\Questions\\[4pt]\S~3.7\\Evaluation\\Plan}
\begin{response}
\underline{Response:} We agree that this is an interesting point and worth investigating, but choosing a single model trained on both feet rather than separate models for the left and right foot was an intentional design decision from the start, for three reasons.

First, splitting the data by laterality would reduce the number of training images from about 600 to about 300 per model, which would mean losing a large amount of data relative to the limited dataset available in this study.

Second, instead of reducing the amount of training data, we added an extra detail to Research Question~1 (RQ1) by comparing a single model trained on images from both feet in their original orientation (S1) with the same model trained after horizontally flipping the left foot to match the orientation of the right foot (S2). The result of RQ1 will provide empirical evidence on whether horizontal flipping helps the model learn better.

Third, although the model in this study makes predictions at the foot level, the predictions from both feet of the same patient can still be used to reflect the prediction at the patient level. For example, if patient 001 is predicted as positive for the right foot but negative for the left foot, we can still conclude that this patient should be referred to the physician for further examination at the diabetic clinic.

For these reasons, we believe that using a single model is sufficient and more appropriate for this study.
\end{response}

% ── C8 ────────────────────────────────────────────────────────────────────────
\cid{C8} The feature counts reported in the baseline feature table appear lower than expected. Descriptors such as LBP and HOG typically produce hundreds or thousands of raw dimensions, yet the table shows only a small number of features per descriptor. It would be helpful to clarify the actual feature dimensions being used and how the dimensionality was reduced.

\mref{\S~3.4.1\\Baseline\\Feature\\Extraction}
\begin{response}
\underline{Response:} We thank the Committee for this observation and agree that the proposal did not explain the feature extraction clearly enough. The raw versions of GLCM, LBP, and HOG are high-dimensional. GLCM is a co-occurrence matrix whose size depends on the number of quantized gray levels, so when the image is quantized to 8 gray levels it becomes an $8\times8$ matrix, while 256 gray levels would yield a $256\times256$ matrix. Because GLCM is also computed at multiple orientations, the number of matrices increases further. LBP produces a histogram of local binary patterns, and with the standard setting of 8 neighboring pixels it can generate up to $2^8 = 256$ possible patterns, so the raw histogram typically contains 256 bins. HOG computes gradient orientation histograms over multiple cells and blocks across the image, which usually results in a vector with several thousand dimensions depending on the image size, cell size, block size, and number of bins.

However, in our baseline method, we do not use these raw high-dimensional representations directly. Following Bansal and Vidyarthi (2024), we summarize the descriptors into compact statistical features. For GLCM, the RGB image is first converted to grayscale and quantized to 8 gray levels, then an $8\times8$ GLCM matrix is computed at four orientations, 0\textdegree, 45\textdegree, 90\textdegree, and 135\textdegree, and four features, Contrast, Correlation, Energy, and Homogeneity, are extracted from each orientation, giving 16 GLCM features in total. For LBP, the image is encoded using the standard neighborhood setting, and the resulting LBP co-occurrence matrix is summarized into eight statistics, namely Mean, Variance, Entropy, Skewness, Kurtosis, Energy, Contrast, and Correlation. For HOG, the image is divided into $8\times8$ cells, gradients are accumulated into orientation histograms, and overlapping $2\times2$ blocks are normalized before concatenation, after which the descriptor is again summarized into eight statistical values. This compact representation is used deliberately to reduce dimensionality and avoid overfitting given the limited dataset size. Table~\ref{tab:baseline_features}, reproduced from Bansal \& Vidyarthi (2024)~[37], summarizes the feature counts for each descriptor and confirms that the compact feature vector is a deliberate design choice of the baseline method, not an error in reporting. A clarifying paragraph with these details has been added to Section~3.4.1.

\begin{center}
\captionsetup{justification=centering, labelfont={color=respblue}, textfont={color=respblue}}
\captionof{table}{Total number of features extracted per descriptor.\\[2pt]
Reproduced from Bansal \& Vidyarthi (2024)~[37].}
\label{tab:baseline_features}
\vspace{0.6em}
\begin{tabular}{lccc}
\toprule
\textbf{Features} & \textbf{Count per Orientation} & \textbf{Orientations} & \textbf{Total} \\
\midrule
GLCM-based & 4 & 4 & 16 \\
LBP-based  & 8 & 1 & 8  \\
HOG-based  & 8 & 1 & 8  \\
\midrule
\textbf{Total} & & & \textbf{32} \\
\bottomrule
\end{tabular}
\par\vspace{0.4em}
\raggedright\small\textsuperscript{[37]}Bansal, N.\ \& Vidyarthi, A.\ (2024). Multivariate Feature-based Analysis of the Diabetic Foot Ulcers Using Machine Learning Classifiers. \textit{Proc.\ 16th Int.\ Conf.\ Contemporary Computing}.
\end{center}
\end{response}

% ── C9 ────────────────────────────────────────────────────────────────────────
\cid{C9} Cross-validation results should be reported as mean~$\pm$~SD across all folds rather than as a single figure, so that the stability of the model's performance can be properly assessed and the results do not rely on any one fold's outcome.

\mref{\S~3.6\\Evaluation Plan\\[4pt]\S~3.7\\Preliminary\\Results}
\begin{response}
\underline{Response:} We agree with this recommendation because, in five-fold cross-validation with a limited dataset size, the performance of each fold can vary depending on which samples are included in that fold. Reporting a single value would hide this variation and could make the model appear more stable than it really is. Reporting the mean~$\pm$~standard deviation across all five folds will make the fold-to-fold variation clearer and provide a better sense of the model's consistency. Accordingly, we have updated the preliminary results section in Section~3.7 and the evaluation plan in Section~3.6.
\end{response}

% ── C10 ───────────────────────────────────────────────────────────────────────
\cid{C10} From what is presented in the proposal, the training approach uses LP-FT, which first freezes the entire backbone and then unfreezes it completely in the second phase. Is this the only strategy being considered? Would it also be possible to unfreeze only a portion of the backbone layers in the second phase, such as 50\% or 30\%?

\mref{\S~1.3\\Research\\Questions\\[4pt]\S~3.6.3.3\\Fine-Tuning\\Strategies\\[4pt]\S~3.7\\Evaluation\\Plan}
\begin{response}
\underline{Response:} We thank the Committee for this question. In the original proposal, we used LP-FT from Dávila et al.\ (2024)~[28] as our starting fine-tuning strategy. In LP-FT, Phase 1 trains only the classification head while the backbone remains frozen, and Phase 2 then unfreezes the entire backbone for full fine-tuning. Partial unfreezing is also a valid strategy, and it is already represented in our study through Gradual Unfreeze Last-to-First (G-LF) and Gradual Unfreeze First-to-Last (G-FL). These strategies unfreeze the backbone gradually rather than all at once, so they let us examine partial unfreezing across different depths of the network, including ranges comparable to the 30 to 50 percent example raised by the Committee.

Following this comment, and in line with Dávila et al.\ (2024)~[28], who found that no single fine-tuning strategy consistently outperformed the others across architectures and medical image types, we expanded the evaluation to include all eight strategies from that study. The comparative results will be reported in RQ1.
\end{response}

{\color{respblue}\raggedright\small\textsuperscript{[28]}Davila, A., Colan, J., \& Hasegawa, Y.\ (2024). Comparison of fine-tuning strategies for transfer learning in medical image classification. \textit{Image and Vision Computing}, 146, 105012.\par}

% ── C11 ───────────────────────────────────────────────────────────────────────
\cid{C11} The proposal mentions the use of McNemar's Test and DeLong's Test as the statistical significance tests for model comparison, but this choice has not been clearly justified. It would be helpful to explain what each test evaluates, why these two specific tests were selected, and why they are appropriate for this study's evaluation design.

\mref{\S~3.7.2\\RQ2: Best CNN\\Model vs.\\Baseline Model}
\begin{response}
\underline{Response:} We thank the committee for this comment. We first clarify what each test evaluates and why both were initially considered. Then, we detail the rationale behind our revised proposal.

Both tests assess whether an observed performance difference between two models is statistically significant, but they operate on fundamentally different quantities. DeLong’s test is a non-parametric method for comparing correlated AUC-ROC values from the same test set. It operates directly on raw probability scores rather than thresholded classifications, computing a Z-statistic from their variance and covariance to assess statistical significance in overall discriminative ability.

In contrast, McNemar’s test evaluates paired binary outcomes at a fixed operating threshold. Using a $2 \times 2$ contingency table, it focuses on discordant pairs, which are instances where only one model predicts correctly. This reveals whether the models significantly differ in their actual decision-making.

Originally, both tests were considered to provide a comprehensive evaluation: DeLong’s test to evaluate the models' overall ranking capabilities, and McNemar’s test to evaluate practical decision-making at a specific threshold.

For the revised proposal, because the proposed model functions as a screening tool where the most critical failure is missing a positive case, we designate Sensitivity as the primary metric for the CNN-versus-Baseline comparison (RQ2). Consequently, we aligned our significance testing with this objective by applying McNemar’s test specifically to the actual positive cases. This directly evaluates whether the two models differ significantly in their sensitivity at the operating threshold. Furthermore, AUC-ROC is now reported alongside other secondary metrics (such as Specificity, PPV, NPV, and F1-Score). Since AUC-ROC is no longer the primary criterion for this specific comparison, applying DeLong’s test is no longer necessary.

\textbf{McNemar's Test for Sensitivity.} On the preliminary INAOE evaluation, restricting McNemar’s test to the 49 positive patients at the default 0.5 threshold and applying the exact binomial form yields the contingency table below.

\begin{center}
\captionsetup{justification=centering, labelfont={color=respblue}, textfont={color=respblue}}
\captionof{table}{Sensitivity McNemar contingency table restricted to the 49 DM patients, Proposed CNN vs.\ Baseline Model.}
\label{tab:mcnemar_prelim}
\vspace{0.6em}
\begin{tabular}{lcc}
\toprule
 & \textbf{Baseline Detects} & \textbf{Baseline Misses} \\
\midrule
\textbf{CNN Detects} & $a = 42$ & $b = 0$ \\
\textbf{CNN Misses}  & $c = 6$  & $d = 1$ \\
\bottomrule
\end{tabular}
\end{center}

The test statistic depends only on the discordant cells $b = 0$ and $c = 6$. The two-sided exact binomial $p$-value is $0.0312$ ($p < 0.05$), so on this preliminary dataset the Baseline achieves significantly higher sensitivity than the proposed CNN. We however acknowledge that this reflects the domain advantage of the thermographic Baseline on the INAOE proxy dataset rather than a fundamental limitation of the CNN, and a more conclusive comparison is anticipated once both models are evaluated on the target podoscope dataset.

We further note that this preliminary comparison is based on a small sample size, with only six discordant DM patients, which limits statistical power, so the result should be read as indicative rather than confirmatory.
\end{response}

\end{document}
